
\section{\BookRecordIndex{T0337}伪循环}\label{sec:6-5}
\zhindex{weixunhuan}{伪循环}

定义\BookRecordIndex{T0143}格罗莫夫--威滕不变量的下一步，是发展一种合适的循环概念。设 $X$ 是光滑
$m$ 维流形。若任意子集 $B\subset X$ 包含在某个映射 $g:W\to X$ 的像中，且流形
$W$ 的每个连通分支维数都不超过 $d$，就称 $B$ 的维数\emph{至多为 $d$}。
\footnote{本书的所有流形都是 $\sigma$-紧的，也就是说，它们都可由可数多个紧集覆盖。}
这时记作 $\dim B\le d$。

\begin{definition}\label{def:6-5-1}
$X$ 中的一个 $d$ 维\emph{\BookRecordIndex{T0337}伪循环}，是定义在定向 $d$ 维流形 $V$ 上的光滑映射
\[
  f:V\longrightarrow X,
\]
使得 $\overline{f(V)}$ 紧，并且
\[
  \dim\Omega_f\le\dim V-2,
  \qquad
  \Omega_f:=
  \bigcap_{\substack{K\subset V\\K\text{ 紧}}}
  \overline{f(V\setminus K)}.
\]
集合 $\Omega_f$ 称为 $f$ 的\emph{\BookRecordIndex{T0001}$\omega$-极限集}。两个 $d$ 维\BookRecordIndex{T0337}伪循环
$f_0:V_0\to X$ 与 $f_1:V_1\to X$ 称为\emph{配边的}，如果存在一个定向
$(d+1)$ 维流形 $W$，满足
\[
  \partial W=V_1\sqcup(-V_0),
\]
以及一个光滑映射 $F:W\to X$，使得 $\overline{F(W)}$ 紧，并且
\[
  F|_{V_0}=f_0,\qquad
  F|_{V_1}=f_1,\qquad
  \dim\Omega_F\le d-1.
\]
\end{definition}
\zhindex{omegajixianji}{$\omega$-极限集}

粗略地说，\BookRecordIndex{T0001}$\omega$-极限集可以理解为 $V$ 的边界在 $f$ 下的像。下一个习题说明，当
$V$ 是一个带边流形 $\overline V$ 的内部，并且 $f$ 连续延拓到 $\overline V$ 时，确实
如此。不过，一般而言，把这一条件直接列入定义会过强。

\begin{exercise}\label{ex:6-5-2}
\begin{enumerate}[label=(\roman*),leftmargin=3em]
\item 证明：$x\in\Omega_f$，当且仅当存在一列在 $V$ 中没有收敛子列的
$v_\nu\in V$，使 $f(v_\nu)\to x$。（这与动力系统中 \BookRecordIndex{T0001}$\omega$-极限集的概念一致。）
证明 $\Omega_f$ 总是紧的。
\item 若 $V$ 是紧带边流形 $\overline V$ 的内部，而 $f$ 连续延拓为
$\overline f:\overline V\to X$，证明
\[
  \Omega_f=\overline f(\partial\overline V).
\]
举例说明，即使 $\dim V\le\dim X$，$\Omega_f$ 仍可能严格大于
$\overline{f(V)}\setminus f(V)$。（当 $d\ge\dim X+2$ 时，每个映射
$f:V\to X$ 都是\BookRecordIndex{T0337}伪循环。）
\sourcefootnote{原书括号内断言只由维数条件自动得到
$\dim\Omega_f\le d-2$，却没有保证定义~\ref{def:6-5-1} 还要求的
$\overline{f(V)}$ 紧；例如恒等映射
$\mathbb R^d\to\mathbb R^d$ 的像闭包并不紧。若另假设 $X$ 紧或 $f(V)$ 预紧，
该断言才成立。原书在例~\ref{ex:6-5-12} 又重复了同一缺条件说法。}
\end{enumerate}
\end{exercise}

$X$ 中\BookRecordIndex{T0337}伪循环的配边类在不交并所给的加法下构成阿贝尔群。零元是在空流形
$V=\varnothing$ 上定义的空映射；$f:V\to X$ 的逆元由反转 $V$ 的定向得到。可以把
$X$ 的\emph{\BookRecordIndex{T0336}伪同调}定义为\BookRecordIndex{T0337}伪循环配边类之群。Schwarz \srcite{359} 与
Zinger \srcite{429} 已经证明，由此得到的同调群与奇异同调一致；不过，我们不会用到
这一事实。

\begin{remark}\label{rem:6-5-3}
为了用\BookRecordIndex{T0337}伪循环 $f:V\to X$ 表示一个 $d$ 维奇异同调类 $\beta$，先用定义在无边界、
$d$ 维有限定向单纯复形 $P$ 上的映射 $f:P\to X$ 表示它。无边界条件意味着，顶维
单形的定向面两两抵消。因此，$P$ 带有一个 $d$ 维基本同调类 $[P]$，并且依定义
\[
  \beta=f_*[P].
\]
现在把 $f$ 逼近为在每个单形上光滑的映射 $f'$。把 $P$ 的 $d$ 维面与
$(d-1)$ 维面之并看成一个光滑 $d$ 维流形 $V$，再把 $f'$ 逼近为跨越
$(d-1)$ 维单形也光滑的映射。可以把这一光滑化 $f''$ 构造成可连续延拓到 $P$。于是
\[
  \Omega_{f''}=f''(P\setminus V),
\]
即 $P$ 中余维至少为二的那些面的像。稍微滥用语言，我们把这类\BookRecordIndex{T0337}伪循环
$f'':V\to X$ 称为\emph{光滑循环}。任意两个表示 $\beta$ 的光滑循环都是配边的。
\end{remark}

当然，若使用有理同调 $H_*(X;\mathbb Q)$，事情会容易一些。Thom 的定理
\srcite{395} 断言：对每个整系数奇异同调类
$\alpha\in H_*(X;\mathbb Z)$，都存在一个整数 $k$，使 $k\alpha$ 可由 $X$ 的某个
紧定向无边子流形表示。因此，$H_*(X;\mathbb Q)$ 有一组基，其中每个元素都可由
$X$ 的光滑闭子流形表示。

\begin{exercise}\label{ex:6-5-4}
设 $X$ 是满足 $\dim X\ge3$ 的光滑流形。证明：每个二维整系数同调类
$A\in H_2(X;\mathbb Z)$ 都可由一个定向嵌入曲面表示。

\emph{提示：}依奇异同调的定义，$A$ 可由定义在无边界紧二维单纯复形上的连续映射
表示。把成对\BookRecordIndex{T0254}粘合的边整理起来，并在必要时分离顶点的不同局部分支，就可用一个无边界
光滑紧二维流形来替代该复形。
\sourcefootnote{原书直接称每个这样的二维复形本身都可赋予闭光滑曲面结构；一般复形的
一个顶点可能有多个圆周状局部分支，须先分离这些分支（即作顶点正规化）。这一处理不改变
所表示的二维循环，并给出正文中的闭曲面。}
所以，$A$ 可由定义在闭二维流形 $\Sigma$ 上的连续映射
$f:\Sigma\to X$ 表示。把 $f$ 逼近为光滑映射，再用一般位置论证使 $f$ 成为浸入；当
$\dim X=4$ 时，它只有有限多个\BookRecordIndex{T0167}横截自交点。当 $\dim X=3$ 或 $4$ 时，用适当改变
$\Sigma$ 的局部手术消去自交；当 $\dim X>4$ 时，用一般位置论证得到嵌入。
\end{exercise}

若两个\BookRecordIndex{T0337}伪循环 $e:U\to X$ 与 $f:V\to X$ 满足
\[
  \Omega_e\cap\overline{f(V)}=\varnothing,
  \qquad
  \overline{e(U)}\cap\Omega_f=\varnothing,
\]
并且
\[
  e(u)=f(v)=x
  \qquad\Longrightarrow\qquad
  T_xX=\operatorname{im}de(u)+\operatorname{im}df(v),
\]
就称它们\emph{\BookRecordIndex{T0281}强横截}。若 $e$ 与 $f$ \BookRecordIndex{T0281}强横截，则
\[
  \{(u,v)\in U\times V\mid e(u)=f(v)\}
\]
是维数为 $\dim U+\dim V-\dim X$ 的紧流形。特别地，若 $U$ 与 $V$ 的维数互补，
这个集合是有限集。
\zhindex{qianghengjie}{强横截}

\begin{lemma}\label{lem:6-5-5}
设 $e:U\to X$ 与 $f:V\to X$ 是维数互补的\BookRecordIndex{T0337}伪循环。
\begin{enumerate}[label=(\roman*),leftmargin=3em]
\item 存在\BookRecordIndex{T0300}剩余集
\[
  \operatorname{Diff}_{\mathrm{reg}}(X,e,f)
  \subset\operatorname{Diff}(X),
\]
使得对每个 $\phi\in\operatorname{Diff}_{\mathrm{reg}}(X,e,f)$，$e$ 都与
$\phi\circ f$ \BookRecordIndex{T0281}强横截。
\item 若 $e$ 与 $f$ \BookRecordIndex{T0281}强横截，则集合
\[
  \{(u,v)\in U\times V\mid e(u)=f(v)\}
\]
有限。这时定义 $e$ 与 $f$ 的\emph{交数}为
\[
  e\cdot f
  =\sum_{\substack{u\in U,\ v\in V\\e(u)=f(v)}}\nu(u,v),
\]
其中 $\nu(u,v)$ 是 $e(U)$ 与 $f(V)$ 在 $e(u)=f(v)$ 处的\BookRecordIndex{T0201}局部交数。
\item 交数 $e\cdot f$ 只依赖于 $e,f$ 的配边类。
\end{enumerate}
\end{lemma}
\zhindex{jiaoshu}{交数}

\begin{proof}
这可由微分拓扑中的标准论证证明，例如见 Milnor \srcite{289} 和
Guillemin--Pollack \srcite{165}。下面列出要点。断言 (i) 可用与
第~\ref{sec:6-2}节相同的技术证明，只须把参数空间 $\mathcal J$ 换成
$\operatorname{Diff}(X)$。先证明映射
\[
  U\times V\times\operatorname{Diff}(X)\longrightarrow X\times X,
  \qquad
  (u,v,\phi)\longmapsto\bigl(e(u),\phi(f(v))\bigr)
\]
横截于对角线；因而普适空间
\[
  \mathcal M
  =\{(u,v,\phi)\mid e(u)=\phi(f(v))\}
  \subset U\times V\times\operatorname{Diff}(X)
\]
是流形。再考察投影 $\mathcal M\to\operatorname{Diff}(X)$ 的\BookRecordIndex{T0409}正则值，即得对一般的
$\phi$，$e$ 与 $\phi\circ f$ 横截。与第~\ref{sec:6-2}节一样，要使这一步严格成立，
须把参数空间完备化为 $C^\ell$ 微分同胚组成的\BookRecordIndex{T0053}巴拿赫流形。

现在选取光滑映射 $e_0:U_0\to X$ 与 $f_0:V_0\to X$，使得
\[
  \dim U_0=\dim U-2,\qquad \Omega_e\subset e_0(U_0),
\]
并对 $f$ 作同样选择。把上述论证分别用于
$(e_0,f)$、$(e,f_0)$ 和 $(e_0,f_0)$，便可\BookRecordIndex{T0326}推出对一般的 $\phi$，
\[
  \overline{e(U)}\cap\phi(\Omega_f)=\varnothing,
  \qquad
  \Omega_e\cap\overline{\phi(f(V))}=\varnothing.
\]
这证明了 (i)。

断言 (ii) 直接来自\BookRecordIndex{T0281}强横截的定义。为证明 (iii)，设\BookRecordIndex{T0337}伪循环 $f:V\to X$ 与空集配边，
相应配边为 $F:W\to X$。与 (i) 的证明一样，可把这个配边选成与 $e$ \BookRecordIndex{T0281}强横截，并且
\[
  \Omega_e\cap\overline{F(W)}=\varnothing,
  \qquad
  \overline{e(U)}\cap\Omega_F=\varnothing.
\]
于是
\[
  Y:=\{(u,v)\in U\times W\mid e(u)=F(v)\}
\]
是紧定向一维流形，其边界为
\[
  \partial Y
  =\{(u,v)\in U\times V\mid e(u)=f(v)\}.
\]
所以 $e\cdot f=0$，引理得证。
\end{proof}

每个 $d$ 维\BookRecordIndex{T0337}伪循环 $e:U\to X$ 按如下方式决定一个同态
\[
  \Phi_e:H_{m-d}(X;\mathbb Z)\longrightarrow\mathbb Z.
\]
依注~\ref{rem:6-5-3}，用光滑循环 $f:V\to X$ 表示同调类
$\beta\in H_{m-d}(X;\mathbb Z)$。任意两个这样的表示都是配边的，所以由
引理~\ref{lem:6-5-5}(iii)，交数
\begin{equation}\label{eq:6-5-1}
  \Phi_e(\beta)=e\cdot f
\end{equation}
与表示 $\beta$ 的循环 $f$ 的选择无关。同样，若 $e_0$ 与 $e_1$ 配边，则对每个循环
$f$ 都有 $e_0\cdot f=e_1\cdot f$。这证明了下述结论。

\begin{lemma}\label{lem:6-5-6}
同态 $\Phi_e$ 只依赖于 $e$ 的配边类。
\end{lemma}

至此已经证明，每个 $d$ 维\BookRecordIndex{T0337}伪循环 $e:U\to X$ 都决定一个同态
$\Phi_e:H_{m-d}(X;\mathbb Z)\to\mathbb Z$，从而决定空间
\[
  H^{m-d}(X)
  :=H^{m-d}(X;\mathbb Z)\big/\text{挠元}
  =\operatorname{Hom}\bigl(H_{m-d}(X;\mathbb Z),\mathbb Z\bigr)
\]
中的一个元素，记作
\[
  a_e:=[\Phi_e]\in H^{m-d}(X).
\]
这定义了一个从\BookRecordIndex{T0337}伪循环配边类到 $H^{m-d}(X)$ 的同态。若 $e$ 实际上是一个光滑循环
（如注~\ref{rem:6-5-3}），则
\[
  a_e=\operatorname{PD}([e]),
\]
也就是说，$e$ 的同调类是 $a_e$ 的庞加莱对偶。一般而言，$e$ 只是同调类
\[
  \alpha_e:=\operatorname{PD}(a_e)\in H_d(X)
\]
的一个\emph{\BookRecordIndex{T0292}弱代表元}。这里称 $d$ 维\BookRecordIndex{T0337}伪循环 $e:U\to X$ 是同调类
$\alpha\in H_d(X)$ 的\BookRecordIndex{T0292}弱代表元，如果对每个同调类
$\beta\in H_{m-d}(X;\mathbb Z)$ 以及每个表示 $\beta$ 的光滑循环 $f$，都有
\[
  e\cdot f=\alpha\cdot\beta.
\]
要使这个定义有意义，还须证明：当 $f:V\to X$ 是按刚才意义弱表示 $\beta$ 的任意\BookRecordIndex{T0337}伪循环
时，公式~\eqref{eq:6-5-1} 仍成立。这个断言并不显然，因为同一同调类
$\beta\in H_d(X)$ 的两个\BookRecordIndex{T0292}弱代表元未必配边。
\footnote{事实上，由上面提到的 Schwarz \srcite{359} 与 Zinger \srcite{429} 的结果，
该结果只有在 $H_d(X)$ 没有挠元时才保证它们配边。}

\begin{lemma}\label{lem:6-5-7}
设 $e:U\to X$ 是 $d$ 维\BookRecordIndex{T0337}伪循环。若 $(m-d)$ 维\BookRecordIndex{T0337}伪循环 $f:V\to X$ 是同调类
$\beta\in H_{m-d}(X)$ 的\BookRecordIndex{T0292}弱代表元，则
\[
  \Phi_e(\beta)=e\cdot f.
\]
\end{lemma}

\begin{proof}
只须证明 $\beta=0$ 的情形。于是设 $f$ 与每个 $d$ 维光滑循环的交数都为零。我们要证明
$f$ 与每个 $d$ 维\BookRecordIndex{T0337}伪循环的交数也为零。为此，由引理~\ref{lem:6-5-5}，可设
$d$ 维\BookRecordIndex{T0337}伪循环 $e:U\to X$ 处于一般位置，即与 $f$ \BookRecordIndex{T0281}强横截。由维数原因，
$\overline{e(U)}$ 不与紧极限集 $\Omega_f$ 相交。取 $\Omega_f$ 的充分小开邻域
$W\subset X$，使其边界光滑且横截于 $f$。于是
\[
  V_0:=f^{-1}(X\setminus W)
\]
是带边紧流形，而 $f$ 在 $V_0$ 上的限制是光滑映射
\[
  f_0:(V_0,\partial V_0)\longrightarrow(X\setminus W,\partial W).
\]
这个映射与 $X\setminus W$ 中每个光滑 $d$ 维循环的交数都为零。因此，存在整数
$\ell>0$，使 $f_0$ 的 $\ell$ 倍在
$H_{m-d}(X\setminus W,\partial W;\mathbb Z)$ 中为边界。这蕴含\BookRecordIndex{T0337}伪循环
$e:U\to X\setminus W$ 与 $f_0$ 的交数为零，从而它与 $f$ 的交数也为零。
\end{proof}

\begin{remark}[模 $2$ 系数]\label{rem:6-5-8}
我们通常使用整系数或有理系数，但有时模 $2$ 计算更方便。上面已经看到，每个 $d$ 维
\BookRecordIndex{T0337}伪循环 $e:U\to X$ 都由交积决定同态
\[
  \Phi_e:H_{m-d}(X;\mathbb Z)\longrightarrow\mathbb Z;
\]
若 $f:V\to X$ 表示同调类 $\beta$，则 $\Phi_e(\beta):=e\cdot f$。模 $2$ 约化给出同态
\[
  \Phi_e:H_{m-d}(X;\mathbb Z)\longrightarrow\mathbb Z/2\mathbb Z,
\]
满足 $\Phi_e(\beta)=e\cdot f\pmod2$。因此，$e$ 唯一决定一个元素
\[
  a_e\in H^{m-d}(X;\mathbb Z/2\mathbb Z),
\]
进而由庞加莱对偶唯一决定 $H_d(X;\mathbb Z/2\mathbb Z)$ 中的元素。读者可以验证：若把
$H_*(X)$ 和 $H^*(X)$ 理解为以 $\mathbb Z/2\mathbb Z$ 为系数的（上）同调，
引理~\ref{lem:6-5-7} 仍成立。
\end{remark}

\paragraph{乘积与\BookRecordIndex{T0166}横截性。}
后面我们将考察乘积流形中的\BookRecordIndex{T0337}伪循环 $\pi\times e:U\to X\times Y$ 与各因子 $X,Y$ 中
\BookRecordIndex{T0337}伪循环之间的关系。本节余下部分讨论一些相关的函子性质。

\begin{lemma}\label{lem:6-5-9}
设 $f:Y\to Z$ 是光滑流形之间的光滑映射。若 $e:U\to Y$ 是 $Y$ 中的\BookRecordIndex{T0337}伪循环，则
$f\circ e:U\to Z$ 是 $Z$ 中的\BookRecordIndex{T0337}伪循环。
\end{lemma}

\begin{proof}
由定义，
\[
  \Omega_{f\circ e}=f(\Omega_e).
\]
这里需要用到 $\overline{e(U)}$ 在 $Y$ 中紧这一事实。
\end{proof}

引理~\ref{lem:6-5-9} 表明，若 $Y$ 是 $Z$ 的闭子流形，则映射 $e:U\to Y$ 是 $Y$ 中的
\BookRecordIndex{T0337}伪循环，当且仅当它与包含映射的复合也是 $Z$ 中的\BookRecordIndex{T0337}伪循环。此外，若
$\pi\times e:U\to X\times Y$ 是\BookRecordIndex{T0337}伪循环，则 $\pi:U\to X$ 和 $e:U\to Y$ 也都是
\BookRecordIndex{T0337}伪循环；但反命题未必成立。例如，若 $U$ 的维数为 $d=4$，而 $X,Y$ 的维数都为 $2$，
则任意光滑映射 $\pi:U\to X$ 和 $e:U\to Y$ 都是\BookRecordIndex{T0337}伪循环，但乘积
$\pi\times e:U\to X\times Y$ 未必是\BookRecordIndex{T0337}伪循环。
\sourcefootnote{原书在双向断言中只假设 $Y$ 是 $Z$ 的子流形；逆向还须假设 $Y$ 在
$Z$ 中闭（等价地，此处的包含映射恰当），才能由像在 $Z$ 中有紧闭包\BookRecordIndex{T0326}推出其在 $Y$ 中
也有紧闭包。例如，$Y=S^2\setminus\{p\}\subset S^2=Z$ 且 $e=\operatorname{id}_Y$
时，复合映射是二维\BookRecordIndex{T0337}伪循环，而 $e:Y\to Y$ 的像闭包并不紧。后文只用正向推论，不受影响。}
\sourcefootnote{这里再次沿用了例~\ref{ex:6-5-12} 前的缺条件说法；要断言任意
$\pi,e$ 都是\BookRecordIndex{T0337}伪循环，仍须假设各自的像有紧闭包（例如 $X,Y$ 紧）。后面的具体反例取
$X=Y=S^2$，不受此缺口影响。}
为处理这种乘积，我们引入下面的\BookRecordIndex{T0166}横截性
概念。

\begin{definition}\label{def:6-5-10}
设 $X$ 是流形，$X_T\subset X$ 是余维为 $\ell$ 的闭子流形。如果存在有限族光滑映射
\[
  g_i:W_i\longrightarrow X,\qquad i=0,\ldots,d-2,
\]
满足下列条件，就称 $d$ 维\BookRecordIndex{T0337}伪循环 $f:U\to X$ 与 $X_T$\emph{\BookRecordIndex{T0294}弱横截}：
\begin{enumerate}[label=(\roman*),leftmargin=3em]
\item 对每个 $i$，$W_i$ 都是 $i$ 维流形，并且
\[
  \Omega_f\subset\bigcup_{i=0}^{d-2}g_i(W_i);
\]
\item $f:U\to X$ 横截于 $X_T$；
\item 当 $d-2-\ell<i\le d-2$ 时，$g_i:W_i\to X$ 横截于 $X_T$。
\end{enumerate}
\end{definition}
\zhindex{ruohengjie}{弱横截}

(ii)、(iii) 中的\BookRecordIndex{T0166}横截性取通常意义。特别地，我们不要求
$\Omega_f\cap X_T=\varnothing$。因此，定义~\ref{def:6-5-10} 的要求弱于
引理~\ref{lem:6-5-5} 前定义的\BookRecordIndex{T0281}强横截性，但强于 (ii) 所表达的通常\BookRecordIndex{T0166}横截性。

\begin{example}\label{ex:6-5-11}
令 $U:=\mathbb C$，$X:=S^2$，并设 $f:U\to S^2$ 是把 $\mathbb C$ 嵌入
$S^2\setminus\{x_\infty\}$ 的球极平面投影嵌入。于是
$\Omega_f=\{x_\infty\}$，所以 $f$ 是\BookRecordIndex{T0337}伪循环。它与
$X_T:=\{x_0\}$ \BookRecordIndex{T0294}弱横截，当且仅当 $x_0\ne x_\infty$；它与嵌入圆周
$X_T\subset S^2$ \BookRecordIndex{T0294}弱横截，当且仅当 $x_\infty\notin X_T$。
\end{example}

\begin{example}\label{ex:6-5-12}
若 $d=\dim U\ge\dim X+2$，则每个光滑映射 $f:U\to X$ 都是\BookRecordIndex{T0337}伪循环。
\sourcefootnote{与习题~\ref{ex:6-5-2}(ii) 的括号说明相同，这一断言还需要
$\overline{f(U)}$ 紧；单有维数不等式并不足够。此处保留原书表述，疑点见前注。}
在这种情形，$f$ 与闭子流形 $X_T\subset X$ \BookRecordIndex{T0294}弱横截，当且仅当它在通常意义下横截于
$X_T$。为看出这一点，当 $i\ne\dim X$ 时取 $W_i=\varnothing$，而当
$i=\dim X\le d-2$ 时取
\[
  W_i:=X,\qquad g_i:=\operatorname{id}_X.
\]
\end{example}

\begin{exercise}\label{ex:6-5-13}
证明：对一般的点 $x\in X$，\BookRecordIndex{T0337}伪循环 $f:U\to X\times Y$ 与
$\{x\}\times Y$ \BookRecordIndex{T0294}弱横截。类似地，可以扰动 $X$ 的任意子流形 $W$，使 $f$ 与
$W\times Y$ \BookRecordIndex{T0294}弱横截。
\end{exercise}

\begin{lemma}\label{lem:6-5-14}
设 $X_T\subset X$ 是余维为 $\ell$ 的闭子流形，而 $f:U\to X$ 是与 $X_T$ \BookRecordIndex{T0294}弱横截的
$d$ 维\BookRecordIndex{T0337}伪循环。令
\[
  U_T:=f^{-1}(X_T)\subset U.
\]
则限制映射 $f_T:U_T\to X$ 是 $X_T$ 中的 $(d-\ell)$ 维\BookRecordIndex{T0337}伪循环，因而也是 $X$ 中的
\BookRecordIndex{T0337}伪循环。
\end{lemma}

\begin{proof}
当 $i\le d-2-\ell$ 时令 $W_{Ti}:=W_i$；当 $i>d-2-\ell$ 时令
$W_{Ti}:=g_i^{-1}(X_T)$。于是对每个 $i$ 都有
$\dim W_{Ti}\le d-\ell-2$，并且
\[
  \Omega_{f_T}\subset\bigcup_i g_i(W_{Ti}).
\]
\end{proof}

\begin{lemma}\label{lem:6-5-15}
设 $X,Y$ 是流形，$X_T\subset X$ 是闭子流形，并且
$\pi\times e:U\to X\times Y$ 是\BookRecordIndex{T0337}伪循环。若 $\pi\times e$ 与
$X_T\times Y$ \BookRecordIndex{T0294}弱横截，则 $\pi:U\to X$ 与 $X_T$ \BookRecordIndex{T0294}弱横截。
\end{lemma}

\begin{proof}
若 $g_i\times h_i:W_i\to X\times Y$ 横截于 $X_T\times Y$，则
$g_i:W_i\to X$ 横截于 $X_T$。
\end{proof}

\begin{example}\label{ex:6-5-16}
引理~\ref{lem:6-5-15} 的逆命题未必成立。例如，令
\[
  U:=\mathbb C^2,\qquad X=Y:=S^2,
\]
设 $f:\mathbb C\to S^2$ 是例~\ref{ex:6-5-11} 中的\BookRecordIndex{T0337}伪循环，并定义
$\pi:U\to S^2$ 与 $e:U\to S^2$ 为
\[
  \pi(z_1,z_2):=f(z_1),\qquad e(z_1,z_2):=f(z_2).
\]
那么 $\pi\times e$ 与 $\pi$ 都是\BookRecordIndex{T0337}伪循环，其极限集分别为
\[
  \Omega_{\pi\times e}
  =\bigl(\{x_\infty\}\times S^2\bigr)
   \cup\bigl(S^2\times\{x_\infty\}\bigr),
  \qquad
  \Omega_\pi=S^2.
\]
所以，$\pi\times e$ 不与 $\{x_\infty\}\times S^2$ \BookRecordIndex{T0294}弱横截；但
$x_\infty$ 不属于 $\pi$ 的像，故由例~\ref{ex:6-5-12}，$\pi$ 与
$X_T:=\{x_\infty\}$ \BookRecordIndex{T0294}弱横截。
\end{example}

\begin{proposition}\label{prop:6-5-17}
设 $X,Y$ 是光滑闭流形，$X_T\subset X$ 是余维为 $\ell$ 的闭子流形。假设
\[
  \pi\times e:U\longrightarrow X\times Y
\]
是与 $X_T\times Y$ \BookRecordIndex{T0294}弱横截的 $d$ 维\BookRecordIndex{T0337}伪循环。那么限制
\[
  e_T:U_T:=\pi^{-1}(X_T)\longrightarrow Y
\]
是 $(d-\ell)$ 维\BookRecordIndex{T0337}伪循环。此外，若 $f:V\to Y$ 是维数满足
\[
  \dim V=\dim Y+\ell-d
\]
的\BookRecordIndex{T0337}伪循环，则
\begin{equation}\label{eq:6-5-2}
  e_T\cdot_Y f
  =(\pi\times e)\cdot_{X\times Y}\bigl([X_T]\times f\bigr).
\end{equation}
\end{proposition}

\begin{proof}
由引理~\ref{lem:6-5-14}，$\pi\times e$ 在 $U_T$ 上的限制
\[
  \pi_T\times e_T:U_T\longrightarrow X\times Y
\]
是 $(d-\ell)$ 维\BookRecordIndex{T0337}伪循环。因此，由引理~\ref{lem:6-5-9}，$e_T$ 是\BookRecordIndex{T0337}伪循环。为证明关于
交数的断言，在 $Y$ 中对 $f$ 作同痕，使它与\BookRecordIndex{T0337}伪循环 $e_T:U_T\to Y$ \BookRecordIndex{T0281}强横截；结论便直接
来自定义。
\end{proof}
